\section{Conclusion}\label{sec:conclusion}

A chokepoint holder's leverage over its patron is bounded by its exit premium, and the patron's demands are bounded by the holder's slack. Because the conditional parts of the patron's threats draw on that one budget, what Europe can obtain for ASML, and what it can be made to give up, turn on the conditional part of its dependence on Washington and on what China can credibly promise.

The result adds a condition to the theory of weaponized interdependence. \citet{farrell2019weaponized} require that a state have jurisdiction over a hub and the domestic institutions to use it. When the hub's holder bargains with a state that depends on the hub, a third condition applies: the hub confers power on its holder only if the holder's exit premium is positive. The Netherlands licenses ASML's exports, so it has the institutions. It has jurisdiction only in part, because Washington reaches the firm through US-origin inputs and US persons (\cref{sec:background}). Whether it meets the third condition depends on how conditional the US guarantee is on chip policy, and that is the quantity our evidence measures least well.

The limitation we would worry about most is information. The model assumes that all three powers know how conditional the guarantee is, and it models the cost of making the guarantee conditional as a reduced-form assurance cost (\cref{prop:conditionality}). If conditionality were private information, Washington would gain from overstating it, its public statements would be partly cheap talk, and European beliefs would set the exit premium. Our results then describe the case in which those beliefs are correct. A model of where the assurance cost comes from (a reputation with other holders, say, or a Congress unwilling to defend allies who help China) could change how much conditionality Washington chooses.

How conditional the guarantee is can be measured. A survey experiment on US foreign-policy elites could describe an ally that has or has not complied with a US export-control demand, hold the ally and the crisis fixed, and ask whether respondents would support defending it. The difference estimates the issue-specific conditionality that the public record does not reveal. The event dataset described in \cref{sec:tests} would test the rationing logic of \cref{cor:rationing} directly, by coding each US demand on a holder of a chipmaking-equipment chokepoint for the instrument threatened, the holder's response, any compensation and any Chinese response.

What the record shows so far fits the jurisdictional half of the argument. What The Hague obtained in December 2024 was a place on Supplement No.~4, which BIS justified on the ground that governments on the list ``have the authority to control key SME items of concern'' and which it can revoke by amending a rule \citep{bis2024fdpr}. What The Hague opposed in May 2026 was the extraterritorial reach of the MATCH Act. Its trade minister told parliament that the cabinet was against the proposal's extraterritorial effect, since ``every country is responsible for its own export control legislation'' \citep[our translation]{tk2026match}.
